\documentclass[10pt]{article}
\usepackage[T1]{fontenc}
\usepackage{lmodern,amsmath,amssymb,amsthm,microtype}
\usepackage[margin=0.65in]{geometry}
\usepackage[colorlinks=true,urlcolor=blue]{hyperref}
\newtheorem{theorem}{Theorem}
\newtheorem{lemma}{Lemma}
\title{An Exact Prime Classification for OEIS A308090}
\author{Qichao Wang\\{\small Hebei University of Technology}\\[0.3em]
Daoyu Dong\\{\small University of Electronic Science and Technology of China}}
\date{October 26, 2026}
\linespread{0.94}
\begin{document}
\maketitle
\begin{abstract}
The factorial gcd sequence in OEIS A308090 equals $n+1$ exactly when $n+1$ is a prime at least 5, and equals 1 otherwise. This proves the recorded converse conjecture and the proposed equality with a prime-indicator sequence, with its stated exceptional index. The proof uses Wilson's theorem and a complete treatment of the composite factorial exception.
\end{abstract}
\section{Complete classification}
Let
\[
 a(n)=\gcd(2^n+n!,3^n+n!,n+1),\qquad n\ge1.
\]
The official entry~\cite{oeis} conjectures that $a(n)=n+1$ implies primality of $n+1$ and observes agreement, except at $n=2$, with A090585.
\begin{lemma}
If $m>4$ is composite, then $m\mid(m-1)!$.
\end{lemma}
\begin{proof}
If $m=ab$ with $1<a<b<m$, both distinct factors occur in the factorial. The remaining case is $m=p^2$ for a prime $p$: if $p$ is odd, the distinct factors $p$ and $2p$ lie below $p^2$, so their product is divisible by $p^2$. The only excluded square is 4.
\end{proof}
\begin{theorem}
For every $n\ge1$,
\[
 a(n)=\begin{cases}n+1,&n+1\text{ is prime and }n\ge4,\\1,&\text{otherwise}.\end{cases}
\]
\end{theorem}
\begin{proof}
Write $m=n+1$. If $m\ge5$ is prime, Wilson's theorem gives $n!\equiv-1\pmod m$, and Fermat's theorem gives $2^n\equiv3^n\equiv1\pmod m$. Both sums are divisible by $m$, so the gcd is $m$.

If $m>4$ is composite, the lemma gives $n!\equiv0\pmod m$. Hence
\[
 a(n)=\gcd(2^n,3^n,m)=1.
\]
Direct evaluation at $m=2,3,4$ gives 1 in each case, completing the proof.
\end{proof}
Thus the conjectured converse holds without exceptions. A090585 has value $n+1$ when $n+1$ is an odd prime and 1 otherwise, so the entries agree except at $n=2$.

The official comment concerning either single sum requires care: divisibility of one sum alone is not a primality criterion. For example, at $n=7$, $8\mid2^7+7!$ although 8 is composite. The theorem concerns the joint gcd and establishes the precise valid statement.
\section{Validation and scope}
The standard-library verifier compares all 99 displayed official terms and independently computes the defining modular gcd for every $1\le n\le5000$. It also checks the supplied 100-row table. Modular arithmetic avoids materializing $n!$ or the exponential powers, but trial-division primality testing is not asserted to have polynomial cost in the bit length of $n$. The proof determines every index, independently of the finite controls. This is a proof of statements recorded in the cited entry, without a claim of literature priority or OEIS acceptance.
\begin{thebibliography}{1}
\bibitem{oeis} OEIS Foundation Inc., A308090, \url{https://oeis.org/A308090}, accessed October 26, 2026.
\end{thebibliography}

\end{document}
